\subsection{Теорема Фредгольма для компактных самосопряженных операторов}

\begin{definition}[Оператор \(A_\lambda\)]
	Пусть
	\[
	A\in\mathcal L(H),
	\qquad
	\lambda\in\mathbb C.
	\]
	Обозначим
	\[
	A_\lambda=A-\lambda I.
	\]
\end{definition}

\begin{lemma}
	Пусть
	\[
	A\in\mathcal K(H)
	\]
	--- самосопряженный оператор и
	\[
	\lambda\neq0.
	\]
	Тогда на подпространстве
	\[
	M=(\ker A_\lambda)^\perp
	\]
	оператор \(A_\lambda\) ограничен снизу:
	\[
	\exists c>0
	\qquad
	\|A_\lambda x\|\geq c\|x\|,
	\qquad x\in M.
	\]
\end{lemma}

\begin{proof}
	\textbf{Хотим получить противоречие с компактностью \(A\).}
	
	Предположим, что \(A_\lambda\) не ограничен снизу на \(M\).
	Тогда существуют
	\[
	x_n\in M,
	\qquad
	\|x_n\|=1,
	\]
	такие что
	\[
	A_\lambda x_n\to0.
	\]
	
	Так как \(A\) компактен, из ограниченной последовательности \(\{x_n\}\)
	можно выбрать подпоследовательность, для которой
	\[
	Ax_{n_k}\to y.
	\]
	
	Но
	\[
	\lambda x_{n_k}
	=
	Ax_{n_k}-A_\lambda x_{n_k}
	\to y.
	\]
	Так как \(\lambda\neq0\),
	\[
	x_{n_k}\to x:=\frac{y}{\lambda}.
	\]
	
	Подпространство \(M\) замкнуто, поэтому
	\[
	x\in M,
	\qquad
	\|x\|=1.
	\]
	
	Переходя к пределу в
	\[
	A_\lambda x_{n_k}\to0,
	\]
	получаем
	\[
	A_\lambda x=0.
	\]
	Следовательно,
	\[
	x\in\ker A_\lambda.
	\]
	
	Таким образом,
	\[
	x\in M\cap\ker A_\lambda=\{0\},
	\]
	что противоречит
	\[
	\|x\|=1.
	\]
\end{proof}

\begin{lemma}
	При тех же условиях
	\[
	\operatorname{Im}A_\lambda
	\]
	замкнут.
\end{lemma}

\begin{proof}
	По предыдущей лемме оператор
	\[
	A_\lambda|_{(\ker A_\lambda)^\perp}
	\]
	ограничен снизу. Так как \((\ker A_\lambda)^\perp\) --- банахово пространство,
	его образ замкнут. Следовательно,
	\[
	\operatorname{Im}A_\lambda
	\]
	замкнут.
\end{proof}

\begin{theorem}[Фредгольма]
	Пусть
	\[
	A\in\mathcal K(H)
	\]
	--- компактный самосопряженный оператор и
	\[
	\lambda\neq0.
	\]
	Тогда
	\[
	H
	=
	\operatorname{Im}A_\lambda
	\oplus_{\perp}
	\ker A_\lambda.
	\]
\end{theorem}

\begin{proof}
	Так как \(A_\lambda\) самосопряжен, имеем стандартное разложение
	\[
	H
	=
	\overline{\operatorname{Im}A_\lambda}
	\oplus_{\perp}
	\ker A_\lambda.
	\]
	
	По предыдущей лемме
	\[
	\operatorname{Im}A_\lambda
	\]
	замкнут, поэтому
	\[
	\overline{\operatorname{Im}A_\lambda}
	=
	\operatorname{Im}A_\lambda.
	\]
	
	Следовательно,
	\[
	\boxed{
		H
		=
		\operatorname{Im}A_\lambda
		\oplus_{\perp}
		\ker A_\lambda.
	}
	\]
\end{proof}